\chapter{Conclusions and Next Steps}
\label{ch:conclusions}

This chapter closes the report by returning explicitly to the research
questions posed in Chapter~\ref{ch:introduction}, stating the activity's
overall scientific conclusion, and setting out what should be done next.
It is a synthesis, not a further results chapter: no new evidence is
introduced here, and every claim below is a restatement, at a bounded
level of confidence, of a conclusion already established in
Chapters~\ref{ch:activities}--\ref{ch:status}.

\section{Answers to the Research Questions}
\label{sec:conclusions-rq}

\textbf{RQ1 (visible-corrosion estimation).} Visible surface corrosion is
the strongest and most consistently learnable image-based quantity in the
available dataset, reached with useful accuracy by three distinct
representations in turn (Chapter~\ref{ch:activities}). This should not be
read as unqualified: a total-rust benchmark in particular carries a
label-adjacency risk that peak-rust benchmarks do not
(Chapter~\ref{ch:results}), and campaign-level visible-corrosion
robustness is supported across both later phases, whereas treatment-level
behaviour is phase-dependent rather than uniformly robust or uniformly
fragile (Chapter~\ref{ch:interpretation}). The conclusion is bounded to
this dataset and these representations, not a claim of universal
generalisation.

\textbf{RQ2 (hidden structural condition).} Structural inference is
substantially weaker and far less transferable than visible-corrosion
estimation. The weaker structural results must be interpreted in light of
the sparse, terminal, and campaign-confounded supervision available for
those targets: structural labels are one measurement per specimen, while
visible-corrosion labels are dense and repeated
(Chapter~\ref{ch:introduction}, Chapter~\ref{ch:dataset}). One
implementation's structural models are mixed-input; a later
implementation's final, robustness-weighted structural selections use
metadata alone; and campaign-holdout evaluation exposes severe transfer
limitations in both (Chapter~\ref{ch:results}). None of this establishes
that no physical relationship exists between surface corrosion and
structural capacity — only that this dataset and design do not isolate one
robustly (Chapter~\ref{ch:interpretation}). The available experiments do
not isolate the relative contribution of data limitations, image
information, and model representation to this weakness.

\textbf{RQ3 (incremental image value for terminal ultimate load).} Under
the pooled setting actually tested, images did not provide reliable
incremental predictive value beyond specimen and exposure metadata
(Chapter~\ref{ch:activities}, Chapter~\ref{ch:results}). A post-onset
analysis showed a limited improvement in some metrics once corrosion had
visibly progressed, but it did not overturn the pooled result, and
neither subgroup nor campaign-transfer evidence was strong enough to
support a general claim that images add value. This conclusion is specific to the terminal
ultimate-load target and the feature representations tested, not a
general statement about image value elsewhere in the project.

\textbf{RQ4 (evaluation robustness).} Which evaluation regime is used
materially changes what can be claimed. A grouped specimen holdout is
useful for estimating generalisation to unseen specimens within the
observed design mixture, but is insufficient on its own to support a
transfer claim. Leave-one-campaign-out is especially revealing for
structural targets, but because campaign identity is entangled with mesh
configuration, chloride level, and exposure duration in this dataset, it
is better understood as an extrapolative stress test than as an ordinary
estimate of future deployment performance
(Chapter~\ref{ch:dataset}, Chapter~\ref{ch:results}). Treatment-level
fragility in visible-corrosion estimation is a genuine finding in one
implementation but is not reproduced in another, and should be treated as
phase-specific rather than a settled, project-wide property
(Chapter~\ref{ch:results}, Chapter~\ref{ch:interpretation}).

\textbf{RQ5 (degradation and proxy-RUL).} The activity produced a working
methodological and diagnostic threshold-screening pipeline, not a
validated remaining-life prediction method. Its structural inputs are
model-derived trajectories, not measurements; no true remaining-useful-life
supervision exists anywhere in the dataset; and the refined pipeline
configuration's own output shows zero \emph{future} threshold crossings —
not zero crossings overall, since many cases are already crossed at
baseline or during observation
(Chapter~\ref{ch:results}, Chapter~\ref{ch:interpretation}). Its
structural inputs are not yet drawn from out-of-fold predictions
(Chapter~\ref{ch:engineering}), which further limits how the resulting
trajectories should be read.

\section{Overall Scientific Conclusion}
\label{sec:conclusions-overall}

The activity began with a broad question: how much can be recovered about
a ferrocement specimen's condition from its repeated corrosion imagery
alone? The evidence gathered across this report answers that question with
a hierarchy of confidence rather than a single number. Visible surface
condition is the quantity most reliably recovered from images. Hidden
structural inference is markedly weaker, a result that is consistent
with — and must be interpreted in light of — its sparse and
campaign-confounded supervision, though the available evidence does not
isolate how much of that weakness is attributable to the supervision
itself as opposed to the image information or model representations used.
Longitudinal degradation and proxy-RUL screening sit
furthest from direct observation and carry correspondingly the least
evidential weight, remaining exploratory and dependent on the model
estimates beneath them (Chapter~\ref{ch:interpretation}).

The central lesson this activity supports is not that one model
outperformed another. It is that how strong a claim can legitimately be
made is jointly determined by how directly the target is observable, how
densely it is supervised, where the label itself came from, whether a
metadata-only baseline was checked before crediting an image, and which
evaluation regime was used to test the result. Read against that lesson,
the project's later decision to refocus the structural research question
away from a sparse, multiply-derived hidden-damage estimate and toward
terminal ultimate load — one of the two structural quantities measured
directly, for every specimen, at a single defined timepoint, and the one
selected for this analysis — was a scientifically sound narrowing of
scope, not a retreat: it exchanged a harder-to-validate question for a
more bounded, directly testable one, and answered it plainly.

\section{Next Steps}
\label{sec:conclusions-next}

Future work falls into three levels, kept explicitly distinct.

\textbf{Immediate research step — prepared, not yet trained.} The
four-class visible-corrosion classification branch has complete,
leakage-safe data preparation but no trained model. The recommended path
is a simple baseline first, with any more complex convolutional or
transformer architecture attempted only once that baseline exists;
validation and test partitions must remain specimen-disjoint and the test
partition untouched until a single final evaluation; and class-sensitive
metrics, not overall accuracy, should be reported given the severity of
the class imbalance (Chapter~\ref{ch:status}).

\textbf{Methodological refinement with existing data — recommended, with
key items still incomplete.} Regenerating degradation inputs from genuinely out-of-fold
structural predictions, checking the physical plausibility of baseline
structural estimates, completing the project's model and feature-artifact
contracts, and repairing the two live reproducibility issues identified in
Chapter~\ref{ch:engineering} are all achievable without new data. None of
this would, on its own, make any remaining-useful-life claim validated;
it would only make the existing exploratory pipeline more transparent and
easier to audit.

\textbf{Data required for stronger scientific claims — a future
requirement, not scheduled work.} The evidence gap this report has
identified is ultimately a data gap, not a modelling gap. Materially
stronger structural or prognostic claims would require additional
independent campaigns, experimental design factors decoupled from
campaign identity, more structural-labelled specimens, repeated structural
measurements over time rather than a single terminal one, and — if
remaining-useful-life is ever to be a directly supervised objective —
observed lifetime or failure-event data. No degree of model refinement
applied to the present dataset can substitute for this missing
supervision.

\section{Closing Statement}
\label{sec:conclusions-closing}

This activity does not deliver a deployment-ready predictor of structural
condition or remaining service life, and it was not positioned to. What it
does deliver is a progressively more rigorously evaluated account of what
repeated corrosion imagery, in this specific experimental setting, can and
cannot support — where the evidence is strong, where it is genuinely
uncertain, and why — together with a progressively more reproducible and
source-traceable basis on which the next modelling and experimental steps
can be built.
